\section{Introduction}
\label{sec:intro}

Many manipulation failures in industrial cells are not failures of geometry but of unobserved physics: a part
is heavier, more slippery or more fragile than the nominal recipe assumes. The quantities involved are
few, have physical units, persist across interactions with the same object or the same lot, and can be probed
at a cost. This suggests a specific architecture: keep what has been learned about each entity in an
explicit \emph{typed} memory (one slot per physical quantity, with uncertainty and provenance), refine it
recurrently as evidence arrives, and let a planner read it when it chooses skill parameters.

This paper is an empirical account of how far that idea gets in simulation, written while the work is still in
progress. It deliberately includes the parts that did not work. The project has gone through one exploratory
sprint and six evaluation rounds; from round~4 on, the test protocol was pre-registered and committed before any
test episode was run.

\paragraph{Questions.}
\begin{enumerate}
  \item Does a typed, structured memory predict better than a generic one of the same size, and is the typed
        slot \emph{causally} used (does overwriting it change predictions and decisions as the simulator says it
        should)?
  \item Does a planner that reads the memory beat scripted baselines in closed loop, in particular a simple
        probe-then-act (\pta{}) script that needs no learning?
  \item Does memory shared across a \emph{lot} of parts change behaviour (fewer probes) and outcomes (more good
        parts)?
  \item Does this transfer to a robot and a task family never seen in training?
  \item Can a pretrained looped language model host the fast physical pathway (\sysI) without disturbing its
        language pathway (\sysII)?
\end{enumerate}

\paragraph{Summary of what we know so far.}
\begin{itemize}
  \item \textbf{Negative:} in the first sprint the typed friction slot collapsed to its prior and was not
        causally used; a physics regulariser made prediction worse; no learned model beat a prior-mean stub in a
        humanoid push loop (Section~\ref{sec:res-sprint}).
  \item \textbf{Positive but narrow:} with a non-bypass architecture and an actor-legal tactile signal, the crush
        posterior of a lot sharpens after one probe in a diagnostic case and is carried to later parts; the typed planner beats the
        generic-memory planner on all four robots, including the held-out one
        (Sections~\ref{sec:res-r7}--\ref{sec:res-r8}).
  \item \textbf{Null against the strongest script:} against \pta{} there is no significant gain in good parts on
        any arm in any round. The single pre-registered win is on the gantry, where a shared call budget makes
        saved probes worth parts (Section~\ref{sec:res-r8}).
  \item \textbf{Ceiling:} absolute success on the arms is low and is limited by the scripted skills and
        controllers, not by knowledge of hidden physics (Section~\ref{sec:res-ceiling}).
  \item \textbf{Pending:} a comparison against ACT-style imitation and compute-matched memory ablations on
        repaired skills, and the closed-loop \sysI/\sysII{} comparison (Section~\ref{sec:res-pending}).
\end{itemize}

Nothing in this paper was run on a physical robot.
